\subsection{From values to behavior}

For a real function we may ask where it is positive, where it increases, what happens near a point, and what happens for very large inputs. For a sequence we ask whether it is bounded, monotone, or approaching a limiting value.

\begin{corebox}[The domain shapes the question]
Real variables can approach a finite point through nearby real values. Sequence indices move through natural numbers and can only tend toward infinity. The common language is functional, but the manner of approaching is different.
\end{corebox}

\begin{examplebox}[Two approaching questions]
The expressions
\[
\lim_{x\to1}f(x)
\qquad\text{and}\qquad
\lim_{n\to\infty}a_n
\]
ask about approached behavior. In the first, a real input approaches a finite point. In the second, a discrete index grows without bound.
\end{examplebox}

\begin{summarybox}[Preparation for limits]
A function assigns outputs to inputs, its domain is part of the object, its graph records input-output pairs, and a sequence is a function on a discrete domain. These ideas provide the language needed for limits and continuity.
\end{summarybox}
